% !TeX root = ../main.tex

\chapter{Background and Related Work}
\label{ch:background}

This chapter provides the necessary background for understanding the methods proposed in this thesis. We review existing approaches to metagenomic taxonomic classification, introduce genomic foundation models, describe parameter-efficient fine-tuning techniques, and discuss attention-based sequence aggregation mechanisms.

\section{Metagenomic Taxonomic Classification}
\label{sec:bg-metagenomics}

Metagenomic taxonomic classification aims to assign each sequencing read from an environmental sample to its source organism's position in the phylogenetic tree. We review three major paradigms for this task.

\subsection{Alignment-Based Methods}

Alignment-based methods classify reads by aligning them to reference genome databases. 
BLAST~\cite{Altschul1990} performs local sequence alignment using heuristic seed-and-extend strategies, achieving high sensitivity but at significant computational cost (hours to days for typical metagenomic datasets). 
DIAMOND~\cite{Buchfink2015} accelerates protein-level alignments through double indexing and optimized scoring, achieving up to 20,000$\times$ speedup over BLAST while maintaining comparable sensitivity. 
MegaBLAST~\cite{Zhang2000} optimizes nucleotide-level alignment for closely related sequences.

The primary limitation of alignment-based methods is their computational cost, which scales poorly with the size of modern metagenomic datasets containing hundreds of millions of reads.

\subsection{$k$-mer-Based Methods}

$k$-mer-based methods represent sequences as sets or distributions of fixed-length subsequences ($k$-mers) and use hash-table lookups for rapid classification.

\textbf{Kraken2}~\cite{Wood2019} constructs a database mapping each $k$-mer ($k=35$) to its lowest common ancestor (LCA) in the taxonomic tree. Classification proceeds by querying each read's $k$-mers against this database and assigning the read to the taxon that covers the most $k$-mers. Kraken2 achieves classification speeds exceeding $10^6$ reads per minute, making it one of the most widely adopted tools in metagenomic analysis.

\textbf{Centrifuge}~\cite{Kim2016} uses a compressed FM-index to perform rapid exact matching of reads against reference genomes, achieving lower memory requirements than Kraken2 while maintaining competitive accuracy.

While $k$-mer methods excel in speed, they rely on exact $k$-mer matches and thus struggle with sequences that diverge from the reference database. They also lack the ability to capture higher-order sequence patterns beyond individual $k$-mers.

\subsection{Deep Learning Methods}

Deep learning methods learn discriminative representations directly from sequence data, potentially capturing complex patterns that elude $k$-mer matching.

\textbf{DeepMicrobes}~\cite{Liang2020} applies a bidirectional LSTM with attention to one-hot encoded nucleotide sequences, achieving competitive performance on species-level classification but requiring substantial training data and compute.

\textbf{MetaTransformer}~\cite{Wichmann2023} is the most directly relevant prior work to this thesis. It employs a custom Transformer architecture that operates on learned 12-mer embeddings. Key design choices include:
\begin{itemize}
    \item \textbf{12-mer tokenization}: Input reads are segmented into non-overlapping 12-mers, each mapped to a learnable embedding vector. The choice of $k=12$ balances vocabulary size ($4^{12} \approx 16.7$M possible 12-mers, though only $\sim$1.4M appear in training data) against sequence granularity.
    \item \textbf{Transformer encoder}: A Transformer encoder (the best-performing configuration uses $d_{\text{model}}=64$ with a single encoder block for $k=13$; the authors also evaluate up to 4 blocks and $d_{\text{model}}=128$) processes the sequence of $k$-mer embeddings.
    \item \textbf{Mean pooling}: Token representations are mean-pooled into a fixed-length vector for classification.
    \item \textbf{Large-scale training}: The model is trained on balanced reads simulated from 2,505 HGR-UMGS genomes using the ART Illumina simulator with a coverage factor of 3 (the exact total number of reads is not reported in the paper; we estimate $\sim$614M based on 300K training steps $\times$ batch size 2,048). The best genus-level model achieves 98.9\% precision and 98.3\% recall on the validation set.
\end{itemize}

A key insight from MetaTransformer is the critical role of training data volume: the authors demonstrate that performance scales substantially with data size, and that balanced sampling across taxa is essential for achieving high macro-level performance.

\section{Genomic Foundation Models}
\label{sec:bg-gfm}

Genomic foundation models (GFMs) are large-scale neural networks pre-trained on vast corpora of DNA sequences using self-supervised learning objectives. Inspired by the success of language models such as BERT~\cite{Devlin2019} and GPT~\cite{Brown2020}, GFMs treat DNA sequences as a ``language'' of nucleotides and learn contextual representations that capture biological structure and function.

\subsection{Nucleotide Transformer}

The Nucleotide Transformer (NT)~\cite{DallaTorre2023} family of models, developed by InstaDeep, are among the largest and most widely evaluated GFMs. This thesis uses the \texttt{nucleotide-transformer-v2-500m-multi-species} variant, which features:

\begin{itemize}
    \item \textbf{Architecture}: An ESM-2-based~\cite{Lin2023} Transformer encoder with 29 layers, 1,024 hidden dimensions, and 16 attention heads.
    \item \textbf{Tokenization}: A 6-mer tokenizer that encodes DNA sequences into non-overlapping 6-nucleotide tokens. This produces a vocabulary of $4^6 = 4,096$ possible tokens plus special tokens.
    \item \textbf{Pre-training}: Masked language modeling (MLM) on 850 genomes from 174 species, spanning $\sim$3.2 billion nucleotides.
    \item \textbf{Total parameters}: 498M parameters.
\end{itemize}

The model produces contextual embeddings for each token in the input sequence, yielding a representation tensor of shape $[\text{batch}, \text{seq\_len}, 1024]$ that captures both local sequence composition and long-range dependencies.

\subsection{Other Genomic Foundation Models}

\textbf{DNABERT}~\cite{Ji2021} was one of the first BERT-style models pre-trained on human genome sequences using overlapping $k$-mer tokenization ($k=3,4,5,6$). While effective for tasks such as promoter prediction, its relatively small scale (110M parameters) and human-genome-only pre-training limit its applicability to diverse metagenomic sequences.

\textbf{DNABERT-2}~\cite{Zhou2024dnabert2} replaces $k$-mer tokenization with Byte Pair Encoding (BPE), eliminating the need to choose a fixed $k$ and enabling more flexible sequence representation. It also adopts a multi-species pre-training strategy.

\textbf{Evo}~\cite{Nguyen2024} is a 7B-parameter model based on the StripedHyena architecture~\cite{Poli2024}, capable of modeling DNA sequences up to 131K nucleotides. While powerful for long-range genomic modeling, its large size makes it challenging to deploy for metagenomic classification tasks involving hundreds of millions of short reads.

\subsection{GFMs for Metagenomic Classification}
\label{sec:bg-gfm-metagenomics}

Despite their success on single-genome tasks, GFMs have seen limited application to metagenomic classification. A na\"ive approach extracts mean-pooled embeddings from a frozen pre-trained model and trains a downstream classifier (e.g., MLP or SVM) on these fixed representations. While simple, this approach discards the token-level structure of the representation, losing positional information that may be discriminative for fine-grained taxonomic distinctions.

This thesis addresses this gap by proposing a token-level classification approach that preserves the full embedding sequence and fine-tunes the backbone with LoRA.

\section{Parameter-Efficient Fine-Tuning}
\label{sec:bg-peft}

Fine-tuning the full parameter set of a 498M-parameter model requires substantial compute and memory, and risks catastrophic forgetting of the pre-trained representations. Parameter-efficient fine-tuning (PEFT) methods address this by adapting only a small subset of parameters while keeping the majority frozen.

\subsection{Low-Rank Adaptation (LoRA)}

LoRA~\cite{Hu2022} is based on the hypothesis that weight updates during fine-tuning lie in a low-rank subspace. For a pre-trained weight matrix $W_0 \in \mathbb{R}^{d \times k}$, LoRA parameterizes the update as:
\begin{equation}
    W = W_0 + \Delta W = W_0 + BA
    \label{eq:lora}
\end{equation}
where $B \in \mathbb{R}^{d \times r}$ and $A \in \mathbb{R}^{r \times k}$ with rank $r \ll \min(d, k)$. During training, only $A$ and $B$ are updated while $W_0$ remains frozen. The forward pass becomes:
\begin{equation}
    h = W_0 x + \frac{\alpha}{r} B A x
\end{equation}
where $\alpha$ is a scaling hyperparameter that controls the magnitude of the adaptation.

Key advantages of LoRA include:
\begin{itemize}
    \item \textbf{Parameter efficiency}: With $r=16$ applied to query, key, and value projections across 29 Transformer layers, only 5.54M parameters (1.11\% of total) are trainable.
    \item \textbf{No additional inference latency}: The adapted weights $W_0 + BA$ can be merged, adding no overhead at inference time.
    \item \textbf{Modular}: Different LoRA adapters can be trained for different tasks and swapped at deployment time.
\end{itemize}

\subsection{Other PEFT Methods}

\textbf{Adapter layers}~\cite{Houlsby2019} insert small trainable bottleneck modules between Transformer layers. While effective, they introduce additional inference latency due to the extra forward-pass computation.

\textbf{Prefix tuning}~\cite{Li2021} prepends learnable ``virtual tokens'' to the input sequence, influencing the attention computation without modifying model weights. It is memory-efficient but can struggle with tasks requiring fine-grained output control.

We choose LoRA for this thesis due to its combination of strong performance, parameter efficiency, and the absence of additional inference latency.

\section{Attention Mechanisms for Sequence Aggregation}
\label{sec:bg-attention}

Given a sequence of token embeddings $\{h_1, h_2, \ldots, h_T\} \in \mathbb{R}^{T \times d}$, a classification task requires aggregating these into a fixed-dimensional representation. The choice of aggregation mechanism directly impacts what information is preserved.

\subsection{Mean Pooling}

The simplest approach averages all token embeddings:
\begin{equation}
    z = \frac{1}{T} \sum_{t=1}^{T} h_t
\end{equation}
Mean pooling is computationally trivial but weights all tokens equally, potentially diluting the contribution of discriminative tokens with non-informative ones.

\subsection{Attention Pooling}

Attention pooling uses learnable query vectors to compute a weighted combination of token embeddings. Given $M$ learnable query vectors $\{q_1, \ldots, q_M\} \in \mathbb{R}^{M \times d_q}$ and linear projections $W_K, W_V \in \mathbb{R}^{d \times d_q}$:
\begin{align}
    K &= H W_K, \quad V = H W_V \\
    \alpha_{m,t} &= \frac{\exp(q_m \cdot k_t / \sqrt{d_q})}{\sum_{t'} \exp(q_m \cdot k_{t'} / \sqrt{d_q})} \\
    z_m &= \sum_{t=1}^{T} \alpha_{m,t} v_t \\
    z &= \text{Concat}(z_1, \ldots, z_M) \in \mathbb{R}^{M \cdot d_q}
\end{align}
This mechanism enables the model to learn which positions in the DNA sequence are most informative for classification, automatically focusing on discriminative $k$-mers while down-weighting uninformative regions.

\subsection{Multi-Head Attention Pooling}

Extending attention pooling with multiple heads allows the model to attend to different aspects of the sequence simultaneously. Each head learns a different query, capturing diverse patterns (e.g., one head may focus on conserved marker genes while another attends to variable regions). The concatenated outputs from all heads provide a rich, multi-faceted representation for downstream classification.

\section{Data Augmentation for DNA Sequences}
\label{sec:bg-augmentation}

\subsection{Reverse-Complement Symmetry in DNA}
\label{sec:bg-rc-symmetry}

DNA is a double-stranded molecule: the two antiparallel strands are linked by Watson--Crick base pairing ($A \leftrightarrow T$, $C \leftrightarrow G$). Given a sequence $s = s_1 s_2 \cdots s_L$ on the forward (5$'$$\to$3$'$) strand, its \emph{reverse complement} $\bar{s} = \overline{s_L} \cdots \overline{s_2} \overline{s_1}$ on the complementary strand encodes identical biological information: both correspond to the same genomic locus and therefore the same organism.

In shotgun metagenomic sequencing, a read is equally likely to originate from either strand. Consequently, a biologically correct classifier must satisfy the \emph{strand symmetry} property:
\begin{equation}
    f(s) = f(\bar{s}) \quad \forall\, s
    \label{eq:strand-symmetry}
\end{equation}
However, standard neural network architectures do not inherently enforce this constraint. Zhou et al.~\cite{Zhou2022rc} systematically demonstrate that deep learning models for genomics frequently violate strand symmetry, producing inconsistent predictions for a sequence and its reverse complement. Ma~\cite{Ma2025rc} further show that even DNA language models fine-tuned with RC data augmentation still exhibit substantial prediction inconsistency.

\subsection{Strategies for Exploiting RC Symmetry}

Three complementary strategies have been proposed to address the strand symmetry problem:

\begin{enumerate}
    \item \textbf{Random RC training augmentation}: During training, each read is replaced with its reverse complement with probability $p=0.5$, exposing the model to both strand orientations. While this encourages---but does not guarantee---strand-invariant representations, it effectively doubles the training data diversity at zero storage cost.

    \item \textbf{RC test-time augmentation (RC TTA)}: At inference, both the forward read $s$ and its reverse complement $\bar{s}$ are processed independently, and their output logits are averaged before taking the argmax prediction:
    \begin{equation}
        \hat{y} = \arg\max_k \bigl[ f_\theta(s) + f_\theta(\bar{s}) \bigr]
        \label{eq:rc-tta-bg}
    \end{equation}
    This ``post-hoc conjoined'' approach~\cite{Zhou2022rc} enforces strand symmetry at inference time without modifying the trained model. It can be understood as a two-view ensemble where both views observe the same biological entity from different strand orientations. By the standard ensemble argument, averaging two correlated but noisy estimates of the same ground-truth label reduces prediction variance and cancels orientation-dependent errors. Zhou et al.~\cite{Zhou2022rc} find that this approach consistently performs well across diverse genomic tasks.

    \item \textbf{RC consistency training}: A regularization loss (e.g., KL divergence) explicitly penalizes divergence between $f_\theta(s)$ and $f_\theta(\bar{s})$ during training, producing a single model that is intrinsically more strand-invariant without requiring double inference at test time~\cite{Ma2025rc}. An alternative is to design \emph{RC-equivariant architectures} that hard-code strand symmetry into the network structure~\cite{Brown2021rc}, though this requires custom layers incompatible with pre-trained foundation model backbones.
\end{enumerate}

In this thesis, we adopt strategies (1) and (2): RC augmentation during training and RC TTA during evaluation. We also experiment with RC consistency training (Section~\ref{sec:exp-techniques}) but find that its interaction with other regularization techniques (e.g., logit adjustment) can be detrimental, and RC TTA alone provides a reliable $+1$--$1.5$ percentage point improvement across all experimental settings (Section~\ref{sec:exp-rc-tta}).

\section{Handling Class Imbalance}
\label{sec:bg-imbalance}

Metagenomic samples exhibit severe class imbalance: a few dominant taxa may constitute the majority of reads while rare taxa contribute only a handful. This imbalance can bias classifiers toward frequent classes, resulting in high micro accuracy but poor macro F1 due to frequent misclassification of rare taxa.

\subsection{Loss-Level Approaches}

\textbf{Class-weighted cross-entropy} scales the loss for each class inversely proportional to its frequency:
\begin{equation}
    \mathcal{L} = -\sum_{k=1}^{K} w_k \, y_k \log \hat{y}_k, \quad w_k \propto 1/n_k
\end{equation}
However, this approach can be unstable with mixed-precision training and large class counts~\cite{Cui2019}.

\textbf{Logit adjustment}~\cite{Menon2021} modifies the output logits by a class-prior-dependent offset:
\begin{equation}
    \tilde{f}_k(x) = f_k(x) + \tau \log \pi_k
    \label{eq:logit-adjustment}
\end{equation}
where $\pi_k$ is the prior probability of class $k$ and $\tau$ is a temperature parameter controlling the strength of adjustment. This approach is theoretically motivated as producing Fisher-consistent estimators under class imbalance.

\subsection{Sampling-Level Approaches}

\textbf{Balanced sampling} constructs each training batch by sampling classes uniformly and then sampling instances within each class. This ensures that rare classes are seen as frequently as common ones during training.

\textbf{Class-aware sampling} (e.g., \texttt{WeightedRandomSampler}) assigns sampling weights inversely proportional to class frequency, achieving a similar effect to balanced batches within standard data loading pipelines.

\section{Summary}

This chapter has reviewed the key concepts underpinning our proposed approach: metagenomic classification methods from alignment-based to deep learning, genomic foundation models that provide rich pre-trained representations, parameter-efficient fine-tuning via LoRA, attention-based sequence aggregation, and strategies for data augmentation and class imbalance. In the next chapter, we describe how these components are integrated into a cohesive classification framework.
